\chapter{反抗爱与救赎}
在《魔圆》的结尾，小圆升入更高层面的存在。
她牺牲的不只是生命，而是自身作为“小圆”的全部存在，以免其他魔法少女最终沦为魔女。
从最终话里焰对丘比的说明，以及《叛逆》中渚和沙耶香的话来看，那些被救下的魔法少女并没有彻底消失，而是以某种形式与小圆一同存在于宇宙之外。
我们无法确知那种存在的形态，却可以确定她们仍有意识，仍能判断和选择。
她们似乎也同时是魔法少女与魔女（这当然也是她们一直以来的真实处境）。

动画本篇，以及《叛逆》开场部分，都把这一变化呈现为积极的发展，算是个圆满结局。
这其实颇令人不解。
沙耶香的整条人物线，几乎都在让她明白：
以牺牲自己来拯救他人，是一种错误。
焰不断试图拯救小圆，结果也只是让两人的处境一次次恶化。
所以，《叛逆》回头追问这种“救赎”本身，并不意外。

影片开头的那段幸福日常，可以看作对两种“小圆的世界”的戏仿：
魔法少女们和谐相处，敌人棘手但终究可以战胜。
“梦魇”作为一种象征人类苦难的新敌人，可以看作是“魔兽”的扭曲变体。
它们像魔兽一样彼此相似，却不再是阴森、细长的人形巨物，而是穿着熊套装、从手臂里射出毛绒玩具的荒诞存在。
在小圆改写后的世界中，击败魔兽可以洁净灵魂宝石。
正因为如此，魔法少女间的合作才终于变得可行，不再像过去那些被魔女支配的时间线那样彼此孤立、难以维系。
而在焰的幻想世界里，击败梦魇后产生的是一片弥散开来的光辉，能够直接净化灵魂宝石。
为了接住更多光，团队行动甚至成了更理想的选择。
更重要的是，小圆创造的魔兽世界第一次展现在我们面前时，伴随的是沙耶香的死亡，以及同伴们对她的哀悼。
可在焰的梦境世界里，沙耶香不仅活着，她的愿望也以某种方式被保留下来：
影片让她与杏子站在一起，或者至少强烈暗示了两人之间的配对关系，从而让仁美和恭介也能够顺利走到一起。
换句话说，焰的幻想世界比小圆创造的世界幸福得多！

与此同时，这个世界也显得更加荒诞可笑。
当然，原因并不只在于梦魇穿着熊玩偶服。
日常景象与诡异气氛的并置，本就是超现实主义艺术常用的表现方式。
魔兽，以及它们出没其间的城市，基本上也属于这一脉络。
例如，无脸或面孔被遮蔽的人，就是马格里特\footnote{%
    勒内·马格里特（1898年11月21日—1967年8月15日），比利时超现实主义画家。
    其作品以诙谐而富于思辨性的符号语言著称，常常挑战观者对现实的先入之见，并对后世插画风格产生了深远影响。
}偏爱的题材之一\cite{gohr2009magritte}。
尽管魔女的视觉形象更加戏谑、更富游戏感，但它们始终与极端暴力紧密相连：
既有魔女对人类和魔法少女的伤害，也有魔法少女对魔女的猎杀。
这种童稚、嬉戏的图像与暴力相互缠绕，构成了一种近乎粗粝而残酷的超现实主义。

然而，梦魇则不同。
它们虽会破坏财物，但战斗结束后，现场并不会留下实际受损的痕迹。
魔法少女们对付梦魇的方式，也不是正面攻击，而是布置陷阱、施加束缚，或用武器将梦魇逼入预设的机关之中。
梦魇被“击败”的方式更像一套既高度仪式化、又没有明确因果依据的程序：
冷开场里由魔法少女筹办的宴席，对付仁美梦魇时关于食物的儿歌式吟唱或游戏，皆是如此。

无害化的暴力、食物、韵律、游戏、随意的仪式，都是无意义文学\footnote{
    译注：Tigges 所论的无意义文学（nonsense literature）指一种故意打破常规、通过奇特组合和语言游戏制造趣味的文学形式。
}的常见要素\cite{tigges1988themes-motifs-nonsense}。
这种文学的核心在于另类逻辑：
在游戏、进餐、礼仪等场景中，规则本身也许是任意的，却必须在内部保持一致，并据此规范人物的行动。
它像梦一样，用一套任意规则替换另一套，并让这些规则的结果以某种合乎逻辑的方式继续展开\footnote{%
    Alison Reike 将这种“范畴错误”列为无意义语言的四种主要类型之一；
    其余三种更直接地属于文字游戏，而“范畴错误”既可发生在语言层面，也可发生在叙事层面。
}\cite{reike1992secret-sense}。
可就在这样一个荒诞世界里，五位魔法少女全都活着，而且过得快乐而安稳。
由此看来，这个看似荒诞的世界，反而比小圆所创造的魔兽世界更加理想。

正如我在间奏二中所说，小圆的净土，也就是她的“来世”或天堂，甚至不及焰的梦境。
因为那是一个没有死亡的世界：没有衰败，没有痛苦，没有腐烂。
沙耶香和渚之所以选择在小圆身边一同具现化，是因为她们渴求的东西只有在衰败与死亡之后才会出现：
对沙耶香而言，是她与杏子的关系；对渚而言，则是奶酪。

焰的梦境还直接构成了对小圆“天堂”的戏仿。
她未经同意便将魔法少女和少数普通人一并拉入自己的世界，将她们困在其中，与宇宙隔绝，却又人为地赋予她们幸福。
她之所以“拯救”她们，是因为她们是小圆所爱的人，焰由此把对小圆的爱延伸到她们身上。
至少对杏子和麻美而言，焰最终相信她们可以杀死此岸的魔女\footnote{%
    译注：即焰的魔女形态，其性质是背德。
    背弃时间，反叛时间，指责时间，审判时间。
}，也开始尊重她们，甚至可能还生出了几分喜欢。

想要拯救某人，就必然想要掌控某人。
沙耶香一旦成为了守护骑士，也就把自己放在审判者的位置上（在面对列车上的两个渣男时，她很可能也成了处刑者）。
焰许愿成为小圆的守护者，结果是一次又一次试图剥夺小圆“成为魔法少女”的选择权。
小圆许愿拯救所有魔法少女，使她们不再沦为自己所对抗的存在，结果则是把自己置于女神的位置。

换言之，成为救世主，必然带有一点暴君性质。
这里说的并不是“帮助”：帮助以被帮助者的同意为前提，而且帮助者暂时处在辅助位置，而非支配位置。
救世主既然是在未获同意的情况下行动，就极有可能误判，正如小圆误判了焰。
看看片头字幕就知道：焰被画成了灰色、近似怪物的身影，跪在地上，而其他魔法少女在一旁翩然起舞。
她无法融入她们的幸福。
她最接近那种幸福的时刻，是电影第一段中那个软弱、羞怯、扎着麻花辫的焰。
即便如此，她也能察觉到某种异样。
等到她再次披散头发，影片剩余部分中的她便再也没有真正幸福过。
原因很简单：漫长到无法叙述的苦难，以及只有她一人记得这些痛苦的事实，已经把她的情感扭曲到近乎无法被拯救的程度。

于是，她转而模仿小圆，把人们夺入自己的结界。
但这与小圆所做的事真的有本质区别吗？
小圆的牺牲是无私之爱，焰的牺牲就是自私吗？
哪一种牺牲更大：牺牲自己的存在，还是牺牲自己的灵魂？
从未存在过更糟，还是成为自己曾经珍视的一切的敌人更糟？

答案当然是：这个问题本身就很荒唐。
价值总是相对的，哪种结局更糟，全取决于视角。
很可能，小圆和焰都认为自己已经献出了所能做出的最大牺牲：
小圆最在意人与人之间的联系，而焰更执着于自己的目标。

可是，当我们在焰的行动中看见小圆行动的扭曲倒影时，小圆的选择也随之显出另一副面孔。
如果一个行动让你得到了自己一直想要的一切，它还能算是真正无私的吗？
小圆得以与她所谓的“大家”同在；
她所爱的人全安然无恙；
她消灭了所有魔女；
她终于成为魔法少女；
她获得了意义，甚至可能比历史上的任何人都更重要。
反观焰，她选择成为“恶魔”，选择把小圆留在这个世界，为此付出的代价恰恰是她唯一珍视的东西：
最终与小圆相守的可能。
从此以后，她们迟早会成为敌人。
那么，难道无私的其实是焰吗？

当然不是。
因为“无私的爱”本来就是一个自相矛盾的说法。
影片之所以展示焰的占有欲，正是为了借它反照出小圆的自私。
爱一个人，就是想保护她，甚至想保护她不受她自己的伤害。
爱一个人，就是想与她共度时间。
爱一个人，就是希望她也渴望你。
如果这种爱以健康的方式表达，并获得对等回应，它当然可以是美好的；
无论是否属于爱情，它都是人与人之间最深的纽带。
但任何纽带都可能反过来成为束缚、控制和支配的工具，虐待关系正由此产生。

纵观整部动画，魔法少女们始终被夹在两种欲望之间：
她们真正想要的东西，以及她们认为自己应该想要的东西。
麻美因为只许愿让自己活下来而非许愿救下父母，痛苦不已。
沙耶香和杏子都为了他人而许愿，却没有许愿让对方理解并珍惜这份帮助，也因此付出了惨重的代价。
丘比之所以选择少女，正在于此：从得到第一个玩偶那一刻起，她们就被教导要照顾别人，被塑造成照料者，被要求为他人的福祉负责\cite{fujimura-fanselow1991good-wives}。
社会早已替丘比完成了前期工作，培养出一群会把社会期待误认作自身愿望的女孩，使她们不再许下自己真正所渴望的东西。
当然，结局并不会因此改变。
仅仅许下愿望，就已经注定了魔法少女要么成为魔女，要么战死。
只是一个错选的愿望会让希望与绝望的循环来得更快罢了。

因此，焰找不到任何方式去许下自己真正的愿望：与小圆在一起。
她只能转而许愿照顾小圆。
第一次，是动画第十话前半段所展示的“原初”时间线结尾；
第二次，则是《叛逆》中她成为“恶魔”的时刻。
两次到最后，她都认为自己除了成为“邪恶”以外别无出路。
由此，《叛逆》合上了这部充满循环的作品中最大的循环：
晓美焰从一个黑暗的、看似反派的、扰乱现状的人物，变成了一个同样黑暗、同样看似反派，却维护现状的人物。
她如今比以往任何时候都更像麻美的黑暗面。
